\chapter{Aprender com IA Sem Atrofiar o Cérebro}
\label{ch:aprender}

\section{IA tutor, IA explicador, IA professor chato}

Usar LLM para estudar pode ser o melhor e o pior que já te aconteceu. Melhor, porque:

\begin{itemize}[leftmargin=*]
\item tens um explicador 24/7 que não se cansa,
\item podes perguntar “porquê?” quantas vezes quiseres sem levar olhos em branco,
\item podes adaptar o nível de explicação ao que realmente sabes (ou achas que sabes).
\end{itemize}

Pior, porque:

\begin{itemize}[leftmargin=*]
\item é tentador pedir “faz o trabalho por mim”,
\item é fácil confundires “ler uma explicação” com “saber fazer”,
\item se o modelo alucinar com confiança e tu não fores capaz de detetar, estás a aprender coisas erradas com entusiasmo.
\end{itemize}

Por isso convém distinguir três modos de uso:

\begin{itemize}[leftmargin=*]
\item \textbf{Tutor}: ajuda-te a planear o estudo, escolher prioridades, construir currículos.
\item \textbf{Explicador}: clarifica conceitos específicos, dá exemplos, cria analogias.
\item \textbf{Professor chato}: faz perguntas, corrige-te, aponta falhas e obriga-te a pensar.
\end{itemize}

Experimenta os três modos e verifica quais ajudam a tua tarefa; não há uma proporção universal de valor.

\begin{keyideabox}{Perceber a explicação não é conseguir aplicar}
Tenta resolver um problema novo sem consultar a resposta. O resultado dá-te informação mais útil do que a sensação de familiaridade ao ler. Desconforto não é uma condição necessária de aprendizagem.
\end{keyideabox}

\section{Níveis de explicação: júnior, sénior e especialista}

Uma das grandes vantagens de LLMs é a capacidade de adaptar o nível de explicação. O mesmo conceito pode ser:

\begin{itemize}[leftmargin=*]
\item explicado a um aluno do secundário,
\item aprofundado para um estudante universitário,
\item discutido em modo “entre dois profissionais”.
\end{itemize}

Se não especificas o nível, o modelo escolhe um meio-termo genérico. Resultado: explicações que nem são suficientemente básicas para iniciantes, nem suficientemente profundas para avançados.

Padrão simples:

\begin{lstlisting}
Explica-me [conceito] em 3 níveis:

Nível 1: Explicação para alguém de 15 anos.
Nível 2: Explicação para um estudante universitário
da área.
Nível 3: Explicação para um profissional que já
conhece o básico, com foco em nuances.
\end{lstlisting}

Depois, podes pedir para expandir apenas o nível que te interessa.

\begin{promptlabbox}{Criar um “menu de níveis” para um tema teu}
\textbf{Objetivo e contexto:} Definir pré-requisitos e progressão de um tema que possas confrontar com materiais de referência.

Escolhe um tema em que te sintas mais ou menos confortável (por exemplo, derivadas, SQL joins, redes TCP/IP).

\medskip

Pede ao modelo:

\begin{lstlisting}
Quero que cries um "menu de explicações" para o tema
[tema]. Para cada nível, define:

a quem se destina,

que pré-requisitos assume,

que tipo de exemplos vai usar.

Depois, dá-me a explicação de Nível 1.
\end{lstlisting}

Se o menu fizer sentido, podes ir desbloqueando os níveis à medida que te sentes pronto. É uma forma simples de estruturar a progressão sem depender do humor do modelo.

\textbf{Resultado esperado e validação:} Verifica correção de cada nível e resolve um exercício sem ajuda antes de avançar. O nível é uma decisão pedagógica, não um diagnóstico automático da tua capacidade.

\textbf{Falhas possíveis:} Usar apenas idade como medida de conhecimento ou avançar por teres reconhecido uma explicação.
\end{promptlabbox}

\section{Construir um plano de estudo 80/20 com IA}

Planeamentos de estudo feitos à mão podem ser ótimos, mas também consomem tempo e energia. A IA pode ajudar a:

\begin{itemize}[leftmargin=*]
\item identificar tópicos prioritários a partir do programa e das dificuldades observadas; «80/20» é uma heurística, não uma proporção demonstrada nesta disciplina,
\item distribuir esses tópicos ao longo de semanas,
\item propor ciclos de revisão e exercícios.
\end{itemize}

Padrão de prompt:

\begin{lstlisting}
Quero um plano de estudo de 6 semanas para a disciplina
de [nome], com foco no exame final.

Dados:

Já sei [lista resumida do que dominas].

Tenho mais dificuldade em [lista].

Posso estudar [n] horas por semana.

Objetivo:

Preparar para uma nota-alvo de [objetivo], sem prometer resultados.

Priorizar o programa e os enunciados que forneco;
nao presumir frequencias de exame desconhecidas.

Formato:

Para cada semana: tópicos, objetivos, atividades
(ler, ver vídeos, fazer exercícios, rever).

Indica, claramente, o que é "núcleo 80/20" e o que
é "extra para nota máxima".
\end{lstlisting}

Depois disso, ajustas:

\begin{itemize}[leftmargin=*]
\item cortas exageros,
\item adaptas ao calendário real (feriados, semanas mais complicadas),
\item acrescentas fontes específicas (livros, apontamentos, enunciados).
\end{itemize}

\begin{checklistbox}{Plano de estudo assistido por IA que não é fantasia}
Antes de aceitares o plano:

\begin{itemize}[leftmargin=*]
\item \textbf{Realismo temporal}: as horas propostas por semana cabem mesmo na tua vida?
\item \textbf{Cobertura de exame}: cruzaste os tópicos com enunciados reais de anos anteriores?
\item \textbf{Margem para revisão}: há tempo reservado para rever, ou é só “stuff novo” até à véspera?
\item \textbf{Prioridades claras}: consegues ver quais são os tópicos críticos vs. “nice to have”?
\end{itemize}
\end{checklistbox}

\section{Quizzes, flashcards e active recall}

Ler explicações é modo passivo. Para aprender a sério, precisas de \emph{active recall}: tentar lembrar e reconstruir ideias sem olhar para a solução.

LLMs são excelentes a gerar:

\begin{itemize}[leftmargin=*]
\item questões de escolha múltipla,
\item perguntas abertas,
\item flashcards (frente/verso),
\item variações de perguntas sobre o mesmo conceito.
\end{itemize}

Padrão prático:

\begin{lstlisting}
A partir dos apontamentos abaixo, cria:

10 flashcards (frente = pergunta curta, verso = resposta
curta, em formato tipo "Q: ... / A: ...").

5 perguntas de desenvolvimento que um professor chato
poderia fazer no exame, com enunciados completos.

[cola apontamentos ou resumo]
\end{lstlisting}

Idealmente, combinas isto com uma rotina:

\begin{itemize}[leftmargin=*]
\item usas uma app de flashcards (Anki, etc.) para repetição espaçada,
\item usas as perguntas de desenvolvimento para simular exame em tempo limitado,
\item só depois lês as soluções da IA.
\end{itemize}

\begin{promptlabbox}{Treino por camadas: lembrar \texorpdfstring{$\to$}{->} explicar \texorpdfstring{$\to$}{->} aplicar}
\textbf{Objetivo e contexto:} Praticar recuperação, explicação e aplicação sobre material de estudo autorizado.

Escolhe um tema concreto. Pede ao modelo:

\begin{enumerate}[leftmargin=*]
\item Uma lista de 10 perguntas rápidas para testar memória (definições, fórmulas, conceitos).
\item 5 perguntas para testar compreensão (explicar com as tuas palavras).
\item 3 problemas para aplicar em situações novas.
\end{enumerate}

Depois, faz o seguinte:

\begin{itemize}[leftmargin=*]
\item responde primeiro sem ajuda,
\item só no fim compara com as respostas do modelo,
\item pede ao modelo para se focar em corrigir apenas onde estiveste mais fraco.
\end{itemize}

\textbf{Resultado esperado e validação:} Responde sem ver as soluções e confirma a chave com o material do curso. Regista erros por conceito para repetir mais tarde.

\textbf{Falhas possíveis:} Ler a solução antes de tentar ou memorizar respostas geradas sem confirmar a sua correção.
\end{promptlabbox}

\section{Simular o professor e o corretor de exames}

Outra aplicação poderosa: usar a IA para simular:

\begin{itemize}[leftmargin=*]
\item um professor que faz perguntas orais difíceis,
\item um corretor a aplicar critérios de avaliação a uma resposta tua.
\end{itemize}

Padrão de “professor oral”:

\begin{lstlisting}
Quero que assumes o papel de professor exigente
na disciplina de [nome].

Vais fazer-me perguntas orais, uma de cada vez.

Começa por perguntas de dificuldade média.

Se eu responder bem, sobe a dificuldade;
se eu responder mal, desce um pouco.

Dá sempre feedback rápido e uma resposta modelo
após cada pergunta.

Começa pela primeira pergunta.
\end{lstlisting}

Padrão de “corretor de exame”:

\begin{lstlisting}
Vou colar uma pergunta de exame e a minha resposta.

Quero que:

expliques como interpretas o enunciado,

dês uma grelha de correção em tópicos,

atribuas uma nota de 0 a 20, justificando,

digas o que eu teria de melhorar para chegar ao 18+.

[pergunta]
[resposta tua]
\end{lstlisting}

\begin{warningboxenv}{O modelo não conhece o professor real}
A IA não sabe exatamente como o teu professor corrige. Mas pode aproximar-se de uma correção “razoável” baseada em critérios gerais. Não uses a nota que a IA te dá como previsão exata; usa como indicador de pontos fracos.
\end{warningboxenv}

\section{Limites e riscos de estudar com IA}

O estudo de Bastani e colaboradores distingue prática assistida e avaliação posterior sem ferramenta, com resultados dependentes do desenho do tutor~\citep{bastani2025}. Não mede uma perda cognitiva universal. Há riscos que o desenho pode procurar reduzir:

\begin{itemize}[leftmargin=*]
\item \textbf{Risco de atalhos}: se usas IA para gerar trabalhos inteiros, estás a sabotar a tua própria aprendizagem. E, em muitos contextos, a violar regras académicas.
\item \textbf{Risco de dependência}: delegar sempre a resolução pode prejudicar a prática das competências que precisas de adquirir; o efeito depende do desenho e do contexto.
\item \textbf{Risco de erro confiante}: a IA pode inventar, e se não tens base para verificar, ficas a memorizar coisas erradas.
\end{itemize}

Contramedidas:

\begin{itemize}[leftmargin=*]
\item estabelece regras pessoais (“só peço solução depois de tentar X minutos sozinho”),
\item usa a IA principalmente como explicador, crítico e criador de exercícios,
\item valida as partes críticas com materiais oficiais e professores.
\end{itemize}

\begin{keyideabox}{IA como ginásio, não como doping}
Usa a IA para treinares melhor: mais variação de exercícios, feedback mais rápido, explicações mais ajustadas. Não a uses para “parecer” em forma sem teres feito o treino.
\end{keyideabox}
